\documentclass{situationist-preprint}

\title{Aristotle's Sea Battle}
\msid{ACT-17-b}
\msfamily{ACT-17}
\msdiscipline{Action, Agency and Free Will}
\mstype{Consequence}
\mskeywords{future contingents; sea battle; De Interpretatione 9; branching time; bivalence; settledness}

\begin{document}
\maketitle

\begin{abstract}
In De Interpretatione 9 Aristotle holds that it is necessary that there will or will not be a sea battle tomorrow, but not necessary that there will be one, nor that there will not. We derive this verdict from the result of ACT-17-a. Today, the sentence that there will be a sea battle tomorrow lacks a truth value, while its disjunction with its negation is true; tomorrow, whichever occurs is settled. Aristotle's thesis that what is, necessarily is when it is becomes moment-relative settledness, and the tenseless truth about the battle, relative to the actual history, gives no support to the fatalist.
\end{abstract}

\section{The text and the derivation}

Aristotle asks whether every affirmation and negation about a particular future event is already true or false, and argues that if it were, nothing would happen by chance or as a result of deliberation \citep[ch.~9, 18b26--19a22]{AristotleDeInt}. He concludes that while it is necessary that there will be or will not be a sea battle tomorrow, it is not necessary that there will be one, nor that there will not, and that of a contradictory pair about such a matter one member need not yet be true or false \citep[19a30--19b4]{AristotleDeInt}. Whether he denied bivalence for future contingents has been disputed \citep{Anscombe1956}; we do not settle the exegetical question, but show that the reading on which he did is a consequence of ACT-17-a.

Let $m$ be a moment today, and let $s$ be the sentence \enquote{there will be a sea battle tomorrow}. Whether there is a battle depends on the decisions of the commanders, which are free choices not yet made at $m$.

\begin{proposition}\label{p:sea}
At $m$: (i) $s$ lacks a truth value; (ii) \enquote{$s$ or not-$s$} is true; (iii) neither $s$ nor not-$s$ is settled.
\end{proposition}

\begin{proof}
By Assumption~2.4 of ACT-17-a, since the battle depends on free choices to be made after $m$, there are histories through $m$ on which there is a battle tomorrow and histories on which there is none. So $s$ is open at $m$ (Definition~2.3 of ACT-17-a), and by Theorem~3.1 of ACT-17-a it lacks a truth value, which is (i); (iii) is the openness itself. By Proposition~3.2 of ACT-17-a, the disjunction is true on every history through $m$, and so true at $m$, which is (ii).
\end{proof}

Proposition~\ref{p:sea} is Aristotle's verdict: the disjunction is necessary, in the sense of settled, and neither disjunct is. Excluded middle holds and bivalence fails, as \citet{Thomason1970} showed the supervaluationist semantics allows, and as \citet{Lukasiewicz1920} did not, since his three-valued logic gives the disjunction no value either.

\begin{proposition}\label{p:when}
At every moment tomorrow after the battle has occurred, or after it has become impossible, $s$ as uttered at $m$ is settled relative to that moment.
\end{proposition}

\begin{proof}
Every history through a later moment $m'$ passes through $m$ and shares $m'$'s past (Definition~2.1 of ACT-17-a). If the battle occurs before $m'$, it occurs on every history through $m'$, so the sentence that it occurs after $m$ is true at $m'$ on every such history; if it has become impossible, it is false on every one.
\end{proof}

This is the moment-relative form of Aristotle's remark that what is, necessarily is when it is \citep[19a23]{AristotleDeInt}. Settledness is indexed to moments: a past and present that are settled relative to each moment, and a future that is open relative to it, as TSM-12 holds of every moment and REL-05-b applies to the fixity of the past.

\section{The tenseless truth}

The tenseless sentence \enquote{a sea battle occurs on day $d$} is true or false relative to each history (Definition~2.2 of ACT-17-a), and in particular relative to the actual one. It may seem to follow that, if the battle actually occurs, $s$ was true at $m$ all along, and the fatalist returns. It does not follow. Relative to the actual history the tenseless sentence is true, but nothing at $m$ singles out the actual history from the others through $m$; a mark that did so would be a thin red line, privileging one history without any ground in what is the case at $m$ \citep{BelnapGreen1994}. The openness of the future is not a state of affairs that neither obtains nor fails to obtain, but the unsettledness at $m$ of the tensed sentence $s$; as MET-19-a also holds. So when we say, after the battle, that the one who predicted it was right, we speak of the tenseless truth relative to the history that has come about, not of a value that $s$ had at $m$.

The battle, finally, is a contingent because it rests on deliberation, which is Aristotle's ground for treating it so \citep[19a7--9]{AristotleDeInt}. Events settled at $m$ by the laws and conditions at $m$, such as tomorrow's tide, are not future contingents, and that their occurrence is settled is not fatalism. Theorem~3.3 of ACT-17-a rejects fatalism about free choices; the sea battle falls under it, and the tide does not.

\section{Conclusion}

Aristotle's sea battle is a direct consequence of the branching account of ACT-17-a. Today the prediction lacks a truth value while the disjunction is true; tomorrow the outcome is settled; and the tenseless truth relative to the actual history gives the fatalist nothing, since no fact today singles that history out.

\begin{thebibliography}{9}\small
\bibitem[Anscombe(1956)]{Anscombe1956} Anscombe, G.~E.~M. (1956). Aristotle and the sea battle. \emph{Mind}, 65(257), 1--15.
\bibitem[Aristotle(De Int.)]{AristotleDeInt} Aristotle. \emph{De Interpretatione}. Trans. J.~L. Ackrill, in \emph{Categories and De Interpretatione}. Clarendon Press, 1963.
\bibitem[Belnap and Green(1994)]{BelnapGreen1994} Belnap, N. and Green, M. (1994). Indeterminism and the thin red line. \emph{Philosophical Perspectives}, 8, 365--388.
\bibitem[{\L}ukasiewicz(1920)]{Lukasiewicz1920} {\L}ukasiewicz, J. (1920). O logice tr\'ojwarto\'sciowej. \emph{Ruch Filozoficzny}, 5, 170--171. English translation, On three-valued logic, in L. Borkowski (ed.), \emph{Selected Works}, North-Holland, 1970.
\bibitem[Thomason(1970)]{Thomason1970} Thomason, R.~H. (1970). Indeterminist time and truth-value gaps. \emph{Theoria}, 36(3), 264--281.
\end{thebibliography}

\end{document}
